% !TeX program = pdflatex
%==============================================================================
%  Lista de Revisão --- Aulas 07 a 10 (Bloco II, classificação)
%  Material de preparação para a Avaliação Teórica 02.
%
%  O gabarito sai de "Lista de revisao 07-10 - gabarito.tex", que apenas liga a
%  opção [gabarito] e inclui este arquivo --- fonte única, dois PDFs.
%
%  O CORPO DE CADA EXERCÍCIO ESTÁ EM exercicios-07-10/ex-NN-slug.tex, como na
%  lista 01-06 (cuja pasta é exercicios/), e é incluído por \input também pela
%  "Avaliacao teorica 02.tex", que sorteou três deles. Assim a prova e a lista
%  não divergem num exercício sorteado; a pontuação por item, que só a prova
%  mostra, vem do comando \pts do estilo.
%
%  Nove exercícios inéditos, dois de cada uma das aulas 07, 08 e 09 e três da
%  aula 10: nenhum repete as listas teóricas e práticas 07 a 10 nem a E3. Usam
%  só o conteúdo dos slides, mais a caixa "o preço do ingênuo", que as notas da
%  aula 07 mantêm de propósito. Cada enunciado é autocontido, para poder ir à
%  prova: não cita aula pelo número, nem outro exercício, e nenhuma solução usa
%  \ref.
%==============================================================================
\providecommand{\gabaritoopt}{}
\documentclass[11pt,a4paper]{article}
\usepackage[\gabaritoopt]{../recursos/latex/estilo-avaliacao}

\begin{document}
\cabecalhorevisao{07 a 10}{Classificadores Gaussianos, Métricas, SVM, KNN e Árvores de
Classificação}

\medskip\noindent
Salvo indicação em contrário, $(\X,Y)$ é um vetor aleatório com $\X\in\R^p$ e $Y$ num
conjunto finito $\mathcal{C}$ de classes; no caso binário, $\mathcal{C}=\{0,1\}$, a classe
$Y=1$ é a positiva e $\eta(\x)=\Prob(Y=1\mid\X=\x)$. A perda é a 0--1, com risco
$\risco(g)=\Prob\big(Y\ne g(\X)\big)$, e a amostra de treino
$\Dados=\{(\X_i,Y_i)\}_{i=1}^n$ é i.i.d. Nos exercícios sobre SVM,
as classes são codificadas por $y\in\{-1,+1\}$.

%==============================================================================
\section*{Aula 07 --- Classificação e Classificadores Gaussianos}
%==============================================================================

\begin{exercicio}
\label{ex:07-logistica}
% Aula 07 --- sob o LDA o log-odds é linear; a recíproca é falsa.
Renata é atuária numa seguradora e quer classificar os segurados de um ramo em dois
perfis, $Y=1$ (risco alto) e $Y=0$ (risco baixo), a partir de uma única covariável: o
tempo $X$, em anos, entre a contratação da apólice e o primeiro sinistro. Ela hesita
entre a análise discriminante linear (LDA) e a regressão logística, dois métodos que
chegam a uma fronteira linear por caminhos diferentes. Ao longo do exercício,
$\sigma(u)=e^u/(1+e^u)$ é a função logística, $\log$ é o logaritmo natural e
$\pi_c=\Prob(Y=c)$.
\begin{itens}
  \item \pts{2,2} Comece pelo caso geral, com $p$ covariáveis: suponha que $\X\in\R^p$ e
        que, como no LDA, $\X\mid Y=c\sim N(\bm\mu_c,\bm\Sigma)$ para $c\in\{0,1\}$, com
        $\bm\Sigma$ simétrica e positiva definida. Mostre que
        $\Prob(Y=1\mid\X=\x)=\sigma(\beta_0+\bbeta^\top\x)$ e escreva $\beta_0$ e
        $\bbeta$ em função de $\pi_0$, $\pi_1$, $\bm\mu_0$, $\bm\mu_1$ e $\bm\Sigma$.
        Conclua que, sob as hipóteses do LDA, o modelo da regressão logística vale
        exatamente. (Escreva $\Prob(Y=1\mid\X=\x)$ pelo Teorema de Bayes, com $f_c$ a
        densidade de $\X$ dado $Y=c$, divida o numerador e o denominador por
        $\pi_1f_1(\x)$ e ponha o resultado na forma $\sigma\big(u(\x)\big)$. Depois,
        substitua as densidades normais em $u(\x)$, expanda as duas formas quadráticas e
        observe o que se cancela porque a covariância é a mesma nas duas classes.)
  \item \pts{1,8} Os tempos de Renata não se parecem com uma normal: são positivos e
        muito assimétricos. Suponha agora que $X\mid Y=c$ tenha distribuição
        exponencial de taxa $\lambda_c$, com densidade $\lambda_c e^{-\lambda_c x}$ para
        $x>0$, com $\lambda_0=1$, $\lambda_1=3$ e $\pi_0=\pi_1=1/2$: em média, o
        primeiro sinistro do perfil de risco alto chega três vezes mais cedo. Mostre que
        $\Prob(Y=1\mid X=x)=\sigma(\log3-2x)$ para $x>0$. O modelo da regressão
        logística vale nesta população, isto é, existem constantes $\beta_0$ e $\beta_1$
        com $\Prob(Y=1\mid X=x)=\sigma(\beta_0+\beta_1x)$ para todo $x>0$? As hipóteses do
        LDA, de que $X$ é normal em cada classe e tem a mesma variância nas duas, valem?
\end{itens}

\begin{solucao}
\textbf{(a)} Pelo Teorema de Bayes, com $f_c$ a densidade de $\X$ dado $Y=c$,
\[
  \Prob(Y=1\mid\X=\x)
  =\frac{\pi_1f_1(\x)}{\pi_0f_0(\x)+\pi_1f_1(\x)}
  =\frac{1}{1+e^{-u(\x)}}=\sigma\big(u(\x)\big),
\]
com
\[
  u(\x)=\log\frac{\pi_1f_1(\x)}{\pi_0f_0(\x)} ,
\]
em que a segunda igualdade sai de dividir numerador e denominador por $\pi_1f_1(\x)>0$.
Basta, portanto, mostrar que o \emph{log-odds} $u(\x)$ é afim em $\x$. As duas densidades
normais têm a mesma constante de normalização, $(2\pi)^{-p/2}(\det\bm\Sigma)^{-1/2}$,
porque a covariância é comum, e ela se cancela no quociente:
\[
  u(\x)=\log\frac{\pi_1}{\pi_0}
        -\tfrac12(\x-\bm\mu_1)^\top\bm\Sigma^{-1}(\x-\bm\mu_1)
        +\tfrac12(\x-\bm\mu_0)^\top\bm\Sigma^{-1}(\x-\bm\mu_0).
\]
Como $\bm\Sigma^{-1}$ é simétrica,
$(\x-\bm\mu)^\top\bm\Sigma^{-1}(\x-\bm\mu)
=\x^\top\bm\Sigma^{-1}\x-2\bm\mu^\top\bm\Sigma^{-1}\x+\bm\mu^\top\bm\Sigma^{-1}\bm\mu$. O
termo quadrático $\x^\top\bm\Sigma^{-1}\x$ aparece nas duas formas e se cancela, e sobra
\[
  u(\x)=\beta_0+\bbeta^\top\x,\qquad
  \bbeta=\bm\Sigma^{-1}(\bm\mu_1-\bm\mu_0),\qquad
  \beta_0=\log\frac{\pi_1}{\pi_0}
     -\tfrac12\big(\bm\mu_1^\top\bm\Sigma^{-1}\bm\mu_1-\bm\mu_0^\top\bm\Sigma^{-1}\bm\mu_0\big).
\]
Logo $\Prob(Y=1\mid\X=\x)=\sigma(\beta_0+\bbeta^\top\x)$, que é exatamente o modelo da
regressão logística. A covariância comum foi usada duas vezes, para cancelar as constantes de
normalização e para cancelar o termo quadrático, e só o segundo uso é indispensável para que
o \emph{log-odds} seja afim. Com covariâncias $\bm\Sigma_0\ne\bm\Sigma_1$, as constantes
deixariam em $u(\x)$ a parcela $\tfrac12\log(\det\bm\Sigma_0/\det\bm\Sigma_1)$, que não
depende de $\x$ e iria para $\beta_0$; já o termo
$-\tfrac12\,\x^\top\big(\bm\Sigma_1^{-1}-\bm\Sigma_0^{-1}\big)\x$ não seria identicamente nulo,
e o \emph{log-odds} seria quadrático em $\x$, como no QDA.

\smallskip
\textbf{(b)} Pela mesma fórmula de Bayes, para $x>0$,
\[
  u(x)=\log\frac{\pi_1\lambda_1e^{-\lambda_1x}}{\pi_0\lambda_0e^{-\lambda_0x}}
      =\log\frac{\pi_1}{\pi_0}+\log\frac{\lambda_1}{\lambda_0}-(\lambda_1-\lambda_0)\,x
      =0+\log3-2x ,
\]
e portanto $\Prob(Y=1\mid X=x)=\sigma(\log3-2x)$. O modelo logístico vale exatamente,
com $\beta_0=\log3\approx1{,}0986$ e $\beta_1=-2$. As hipóteses do LDA falham duas vezes:
$X\mid Y=c$ não é normal (vive em $(0,\infty)$ e é assimétrica), e as variâncias das duas
classes, $1$ e $1/9$, são diferentes. O item (a) mostra que as hipóteses do LDA bastam
para o modelo logístico, e este item mostra que não são necessárias: a logística supõe
apenas que o \emph{log-odds} seja afim, e isso acontece aqui sem nada de normal.
\end{solucao}

\end{exercicio}

\begin{exercicio}
\label{ex:07-ingenuo}
% Aula 07 --- uma covariável duplicada: o Bayes ingênuo conta a mesma evidência duas vezes.
Lucas é analista de crédito e quer prever a inadimplência, $Y=1$, a partir da renda
mensal $X_1$ do cliente, em milhares de reais. Na população de clientes,
$X_1\mid Y=1\sim N(3,1)$ e $X_1\mid Y=0\sim N(5,1)$. Ao juntar duas bases, Lucas fica,
sem perceber, com a renda duas vezes: em $X_1$, em milhares de reais, e em $X_2=1000\,X_1$,
em reais. Ele ajusta um Bayes ingênuo gaussiano às duas colunas. Para separar o erro do
modelo do erro de estimação, suponha que a amostra é tão grande que o ingênuo usa as
distribuições verdadeiras de $X_1$ e de $X_2$ em cada classe. Seja
\[
  L(x_1)=\log\frac{f_1(x_1)}{f_0(x_1)}
\]
a razão de log-verossimilhanças da renda, em que $f_c$ é a densidade de $X_1$ dado $Y=c$
e $\log$ é o logaritmo natural, e sejam $\pi_c=\Prob(Y=c)$ e $\sigma(u)=e^u/(1+e^u)$.
\begin{itens}
  \item \pts{1,0} Mostre que $L(x_1)=8-2x_1$. Mostre que, num cliente com renda $x_1$, o
        \emph{log-odds} verdadeiro é $\log(\pi_1/\pi_0)+L(x_1)$ e o calculado pelo
        ingênuo é $\log(\pi_1/\pi_0)+2L(x_1)$.
  \item \pts{1,0} Suponha $\pi_0=\pi_1=1/2$ e seja $\eta=\Prob(Y=1\mid X_1=x_1)$. Mostre
        que a probabilidade calculada pelo ingênuo é
        $\widehat P=\eta^2/\big(\eta^2+(1-\eta)^2\big)$. Conclua que ela ordena os
        clientes como $\eta$, que classifica cada um como o classificador de Bayes, e que
        $\abs{\widehat P-1/2}>\abs{\eta-1/2}$ sempre que $\eta\notin\{0,1/2,1\}$.
        Calcule $\widehat P$ para $\eta=0{,}6$ e para $\eta=0{,}8$.
  \item \pts{1,2} Na carteira de Lucas, porém, só $20\%$ dos clientes são inadimplentes:
        $\pi_1=0{,}2$. Determine as rendas que o classificador de Bayes classifica como
        inadimplentes e as que o ingênuo classifica assim, e mostre que as duas decisões
        diferem num intervalo de rendas. Calcule a probabilidade de um cliente cair nesse
        intervalo e diga qual das duas decisões erra mais nele. \emph{Dados:}
        $\log2\approx0{,}6931$, $\Phi(-1{,}3466)\approx0{,}0891$,
        $\Phi(-1{,}6931)\approx0{,}0452$, $\Phi(0{,}6534)\approx0{,}7433$ e
        $\Phi(0{,}3069)\approx0{,}6205$, em que $\Phi$ é a função de distribuição da
        normal padrão.
  \item \pts{0,8} Se a renda aparecesse em $m$ colunas, todas proporcionais a $X_1$, o que
        aconteceria com o \emph{log-odds} do ingênuo? E com a sua decisão quando
        $m\to\infty$? Resuma o que o exercício mostra sobre o Bayes ingênuo com
        covariáveis redundantes: o que ele preserva e o que ele estraga.
\end{itens}

\begin{solucao}
\textbf{(a)} Com as densidades $N(3,1)$ e $N(5,1)$, as constantes de normalização se
cancelam no quociente, e
\[
  L(x_1)=-\tfrac12(x_1-3)^2+\tfrac12(x_1-5)^2
        =\tfrac12\big[(x_1-5)^2-(x_1-3)^2\big]
        =\tfrac12\,(-2)(2x_1-8)=8-2x_1 ,
\]
em que se usou $a^2-b^2=(a-b)(a+b)$ com $a=x_1-5$ e $b=x_1-3$.

O par $(X_1,X_2)$ carrega a mesma informação que $X_1$, porque $X_2=1000\,X_1$ é uma
função de $X_1$, e $X_1$ é função do par. Saber que $(X_1,X_2)=(x_1,1000\,x_1)$ é,
portanto, o mesmo que saber que $X_1=x_1$, e o \emph{log-odds} verdadeiro é o da renda
sozinha:
\[
  \log\frac{\Prob(Y=1\mid X_1=x_1)}{\Prob(Y=0\mid X_1=x_1)}
  =\log\frac{\pi_1f_1(x_1)}{\pi_0f_0(x_1)}=\log\frac{\pi_1}{\pi_0}+L(x_1).
\]
O ingênuo, por sua vez, trata as duas colunas como independentes dentro de cada classe e
usa o produto das densidades marginais. A densidade de $X_2=1000\,X_1$ dado $Y=c$ é
$h_c(x_2)=\frac{1}{1000}f_c\big(x_2/1000\big)$, pela fórmula da mudança de variável, e o
fator $1/1000$ se cancela no quociente: $h_1(1000\,x_1)/h_0(1000\,x_1)=f_1(x_1)/f_0(x_1)$. O
\emph{log-odds} do ingênuo é então
\[
  \log\frac{\pi_1f_1(x_1)\,h_1(1000\,x_1)}{\pi_0f_0(x_1)\,h_0(1000\,x_1)}
  =\log\frac{\pi_1}{\pi_0}+2L(x_1).
\]
A mudança de unidade não altera a razão de verossimilhanças, e a segunda coluna repete a
evidência da primeira; o ingênuo a soma de novo, como se fosse uma testemunha
independente.

\smallskip
\textbf{(b)} Com prioris iguais, $\log(\pi_1/\pi_0)=0$, o \emph{log-odds} verdadeiro é
$u=L(x_1)$ e o do ingênuo é $2u$. Como $\eta=\sigma(u)$, vale $e^u=\eta/(1-\eta)$, e
\[
  \widehat P=\sigma(2u)=\frac{e^{2u}}{1+e^{2u}}
  =\frac{\eta^2/(1-\eta)^2}{1+\eta^2/(1-\eta)^2}
  =\frac{\eta^2}{\eta^2+(1-\eta)^2}.
\]
A função $u\mapsto\sigma(2u)$ é estritamente crescente, e $u\mapsto\sigma(u)$ também; logo
$\widehat P$ é uma função estritamente crescente de $\eta$, e os dois ordenam os clientes
da mesma maneira. Em particular, $\widehat P>1/2\iff 2u>0\iff u>0\iff\eta>1/2$, e o
ingênuo classifica cada cliente como o classificador de Bayes.

Para $1/2<\eta<1$, dividindo por $\eta>0$,
\[
  \widehat P>\eta
  \iff \eta>\eta^2+(1-\eta)^2
  \iff 2\eta^2-3\eta+1<0
  \iff (2\eta-1)(\eta-1)<0,
\]
o que vale, porque o primeiro fator é positivo e o segundo é negativo. Para
$0<\eta<1/2$, a simetria $\widehat P(1-\eta)=1-\widehat P(\eta)$ dá $\widehat P<\eta$. Nos
dois casos $\widehat P$ fica mais longe de $1/2$ que $\eta$. Com $\eta=0{,}6$,
$\widehat P=0{,}36/0{,}52\approx0{,}692$; com $\eta=0{,}8$,
$\widehat P=0{,}64/0{,}68\approx0{,}941$. O ingênuo acerta a ordem e a decisão, mas exagera
a certeza.

\smallskip
\textbf{(c)} Com $\pi_1=0{,}2$, $\log(\pi_1/\pi_0)=\log(1/4)=-2\log2$. O classificador de
Bayes declara inadimplente quando o \emph{log-odds} verdadeiro é positivo:
\[
  -2\log2+8-2x_1>0\iff x_1<4-\log2\approx3{,}3069 .
\]
O ingênuo declara quando o seu \emph{log-odds} é positivo:
\[
  -2\log2+2(8-2x_1)>0\iff x_1<4-\tfrac12\log2\approx3{,}6534 .
\]
As decisões diferem no intervalo $I=[3{,}3069;\ 3{,}6534)$, de rendas entre cerca de
R\$~3\,307 e R\$~3\,653: ali o ingênuo declara inadimplente quem o classificador de Bayes
declara adimplente. Com prioris iguais isso não acontecia, porque multiplicar o
\emph{log-odds} por $2$ não muda o seu sinal; com prioris diferentes, o termo
$\log(\pi_1/\pi_0)$ não é dobrado, e o ponto em que a soma troca de sinal se desloca.

Pela lei da probabilidade total, com os dados do enunciado,
\begin{align*}
  \Prob(X_1\in I\mid Y=0)&=\Phi(3{,}6534-5)-\Phi(3{,}3069-5)\approx0{,}0891-0{,}0452=0{,}0439,\\
  \Prob(X_1\in I\mid Y=1)&=\Phi(3{,}6534-3)-\Phi(3{,}3069-3)\approx0{,}7433-0{,}6205=0{,}1228,
\end{align*}
e $\Prob(X_1\in I)\approx0{,}8\cdot0{,}0439+0{,}2\cdot0{,}1228\approx0{,}035+0{,}025=0{,}060$:
a decisão muda para cerca de $6\%$ dos clientes. Em $I$, o ingênuo erra nos adimplentes,
com probabilidade $\Prob(Y=0,\,X_1\in I)\approx0{,}035$, e o classificador de Bayes erra nos
inadimplentes, com probabilidade $\Prob(Y=1,\,X_1\in I)\approx0{,}025$. O ingênuo erra mais,
e a sua probabilidade de erro fica cerca de $0{,}011$ acima da de Bayes, como tinha de ser,
porque nenhum classificador erra menos que o de Bayes.

\smallskip
\textbf{(d)} Com $m$ colunas proporcionais a $X_1$, o argumento do item (a) se repete para
cada uma, e o \emph{log-odds} do ingênuo passa a $\log(\pi_1/\pi_0)+m\,L(x_1)$. Quando
$m\to\infty$, o segundo termo domina sempre que $L(x_1)\ne0$: o \emph{log-odds} vai a
$+\infty$ se $L(x_1)>0$ e a $-\infty$ se $L(x_1)<0$, e a probabilidade calculada vai a $1$ ou
a $0$. A decisão passa a ser $x_1<4$, a de prioris iguais, como se a priori tivesse sumido.

O Bayes ingênuo com covariáveis redundantes preserva a ordenação das observações, porque o
seu \emph{log-odds} continua sendo uma função crescente do verdadeiro. Estraga as
probabilidades, que saem empurradas para $0$ e $1$, e, quando as classes têm tamanhos
diferentes, estraga também algumas decisões, porque a evidência é contada mais de uma vez e
a priori, uma vez só. Covariáveis apenas correlacionadas produzem uma versão mais branda do
mesmo efeito; com colunas repetidas ou quase repetidas, convém retirá-las antes do ajuste ou
usar um modelo que leve a correlação em conta.
\end{solucao}

\end{exercicio}

%==============================================================================
\section*{Aula 08 --- Métricas para Classificação}
%==============================================================================

\begin{exercicio}
\label{ex:08-sorteio}
% Aula 08 --- classificadores sorteados: combinações convexas de pontos ROC.
Diego é analista antifraude num banco e tem dois sistemas que já foram avaliados, $A$ e
$B$. Tome a fraude como a classe positiva, $Y=1$. As taxas de falsos positivos e de
verdadeiros positivos são $(\mathrm{FPR},\mathrm{TPR})=(0{,}02;\ 0{,}50)$ para $A$ e
$(0{,}08;\ 0{,}80)$ para $B$. O banco aceita mandar para revisão manual no máximo $5\%$ das
transações legítimas, e nenhum dos dois sistemas aproveita esse limite: $A$ fica abaixo
dele e $B$ o ultrapassa. Diego então propõe sortear: para cada transação, uma moeda com
probabilidade $\lambda$ de dar cara, lançada independentemente de tudo o mais, escolhe se a
decisão será a de $A$ (cara) ou a de $B$ (coroa). Seja $g_\lambda$ esse classificador.
\begin{itens}
  \item \pts{1,6} Mostre que
        $\mathrm{FPR}(g_\lambda)=\lambda\,\mathrm{FPR}(A)+(1-\lambda)\,\mathrm{FPR}(B)$ e
        que o mesmo vale para a TPR, isto é, que o ponto de $g_\lambda$ no plano
        (FPR, TPR) está no segmento que liga os pontos de $A$ e de $B$. Determine $\lambda$
        para que a taxa de falsos positivos seja exatamente $0{,}05$ e calcule a taxa de
        verdadeiros positivos obtida. Que classificadores sorteados dão os pontos da
        diagonal?
  \item \pts{1,2} Um terceiro sistema, $C$, tem $(\mathrm{FPR},\mathrm{TPR})=(0{,}70;\ 0{,}40)$.
        Mostre que o classificador que inverte todas as decisões de $C$ tem o ponto
        $(0{,}30;\ 0{,}60)$ e, em geral, que inverter as decisões leva
        $(\mathrm{FPR},\mathrm{TPR})$ a $(1-\mathrm{FPR},1-\mathrm{TPR})$. Se um escore tem
        curva ROC abaixo da diagonal, com área $\mathrm{AUC}<1/2$, qual é a curva e qual é a
        área do escore com o sinal trocado?
  \item \pts{1,2} Com custos $\ell_0$ por falso negativo e $\ell_1$ por falso positivo, o
        custo esperado de um classificador $g$ é
        $\risco_\ell(g)=\ell_1\Prob(g(\X)=1,\,Y=0)+\ell_0\Prob(g(\X)=0,\,Y=1)$. Mostre que
        $\risco_\ell(g_\lambda)=\lambda\risco_\ell(A)+(1-\lambda)\risco_\ell(B)$ e conclua
        que, sem a restrição do banco, sortear nunca custa menos que o melhor dos dois
        sistemas. Para que serve, então, o sorteio?
\end{itens}

\begin{solucao}
\textbf{(a)} Seja $Z$ a moeda, com $\Prob(Z=1)=\lambda$, independente de $(\X,Y)$. Pela lei
da probabilidade total, condicionando em $Z$,
\begin{align*}
  \Prob\big(g_\lambda(\X)=1\mid Y=0\big)
  &=\lambda\,\Prob\big(g_A(\X)=1\mid Y=0,\,Z=1\big)\\
  &\quad+(1-\lambda)\,\Prob\big(g_B(\X)=1\mid Y=0,\,Z=0\big),
\end{align*}
em que se usou $\Prob(Z=1\mid Y=0)=\Prob(Z=1)=\lambda$, pela independência. Também pela
independência, condicionar em $Z$ não altera a distribuição de $\X$ dado $Y$, e as duas
probabilidades condicionais são as taxas de falsos positivos de $A$ e de $B$. Assim,
$\mathrm{FPR}(g_\lambda)=\lambda\,\mathrm{FPR}(A)+(1-\lambda)\,\mathrm{FPR}(B)$, e a mesma
conta, condicionando em $Y=1$, dá a TPR. O ponto de $g_\lambda$ é a combinação convexa dos
pontos de $A$ e de $B$ com pesos $\lambda$ e $1-\lambda$, e percorre o segmento entre eles
quando $\lambda$ vai de $0$ a $1$.

Para a taxa de falsos positivos valer $0{,}05$,
\[
  0{,}02\,\lambda+0{,}08\,(1-\lambda)=0{,}05\iff 0{,}06\,\lambda=0{,}03\iff\lambda=\tfrac12 ,
\]
e a taxa de verdadeiros positivos é $\tfrac12\cdot0{,}50+\tfrac12\cdot0{,}80=0{,}65$: Diego usa
o limite do banco inteiro e detecta $65\%$ das fraudes, contra $50\%$ de $A$. A diagonal
vem do sorteio entre os dois classificadores triviais, $g\equiv0$, no ponto $(0,0)$, e
$g\equiv1$, no ponto $(1,1)$: sortear $g=1$ com probabilidade $q$ dá o ponto $(q,q)$.

\smallskip
\textbf{(b)} Se $g'=1-g$, então $g'(\X)=1$ exatamente quando $g(\X)=0$, e
\[
  \mathrm{TPR}(g')=\Prob\big(g(\X)=0\mid Y=1\big)=1-\mathrm{TPR}(g),
  \qquad
  \mathrm{FPR}(g')=1-\mathrm{FPR}(g).
\]
Para $C$, o ponto invertido é $(1-0{,}70;\ 1-0{,}40)=(0{,}30;\ 0{,}60)$, acima da diagonal: um
sistema que erra sistematicamente carrega informação, desde que se faça o contrário do que
ele diz. A aplicação $(a,b)\mapsto(1-a,1-b)$ é a simetria em relação ao ponto
$(\tfrac12,\tfrac12)$. Trocar o sinal do escore troca ``escore acima do corte'' por
``escore abaixo do corte'', o que inverte as decisões em cada corte; a curva do escore com o
sinal trocado é, portanto, a simétrica da original em relação a $(\tfrac12,\tfrac12)$. Essa
simetria leva a região do quadrado unitário abaixo de uma curva na região acima da curva
simétrica e preserva áreas. A área sob a curva nova é, então, a área acima da curva antiga,
$1-\mathrm{AUC}>1/2$.

\smallskip
\textbf{(c)} Condicionando de novo na moeda, e usando a independência,
\[
  \Prob\big(g_\lambda(\X)=1,\,Y=0\big)
  =\lambda\,\Prob\big(g_A(\X)=1,\,Y=0\big)+(1-\lambda)\,\Prob\big(g_B(\X)=1,\,Y=0\big),
\]
e o mesmo para $\Prob(g_\lambda(\X)=0,\,Y=1)$. Multiplicando a primeira igualdade por
$\ell_1$, a segunda por $\ell_0$ e somando,
$\risco_\ell(g_\lambda)=\lambda\risco_\ell(A)+(1-\lambda)\risco_\ell(B)$. Uma média ponderada
de dois números nunca é menor que o menor deles, e o sorteio nunca custa menos que o melhor
dos dois sistemas: sem restrição, basta escolher um deles. O sorteio serve para atingir um
ponto que nenhum dos dois atinge, como a taxa de falsos positivos de exatamente $0{,}05$
permitida pelo banco: respeitando esse limite, $g_{1/2}$ detecta $65\%$ das fraudes, e $A$,
só $50\%$.
\end{solucao}

\end{exercicio}

\begin{exercicio}
\label{ex:08-exame}
% Aula 08 --- doença rara: o classificador de Bayes é o trivial, e os custos mudam isso.
Cecília é médica e coordena o rastreamento de uma doença que atinge $1\%$ da população:
$\Prob(Y=1)=0{,}01$, com $Y=1$ para doente. O exame disponível tem resultado $X\in\{0,1\}$ e,
seguindo o resultado (isto é, declarando doente quem tem $X=1$), detecta $90\%$ dos doentes
e dá alarme falso para $10\%$ dos saudáveis:
$\Prob(X=1\mid Y=1)=0{,}9$ e $\Prob(X=1\mid Y=0)=0{,}1$. Seja
$\eta(x)=\Prob(Y=1\mid X=x)$.
\begin{itens}
  \item \pts{1,2} Calcule $\eta(1)$ e $\eta(0)$.
  \item \pts{1,2} Determine o classificador de Bayes, o que minimiza a probabilidade de erro.
        Compare a probabilidade de erro dele com a do classificador que segue o exame. Um
        colega conclui que o exame não serve para nada. O que há de errado nessa conclusão?
  \item \pts{1,6} Sejam $\ell_0$ o custo de deixar passar um doente e $\ell_1$ o de um alarme
        falso, e $r=\ell_0/\ell_1$. Lembre que o classificador que declara doente quando
        $\eta(x)>\ell_1/(\ell_0+\ell_1)$ minimiza o custo esperado. Mostre que esse
        classificador segue o exame, isto é, declara doente exatamente quem tem $X=1$, se e
        somente se $11<r\le891$, e diga o que ele faz fora dessa faixa. Para $\ell_0=20$ e
        $\ell_1=1$, calcule o custo esperado de declarar todos saudáveis, de seguir o exame e
        de declarar todos doentes, numa população de $10\,000$ pessoas.
\end{itens}

\begin{solucao}
\textbf{(a)} Pelo Teorema de Bayes,
\[
  \eta(1)=\frac{0{,}9\cdot0{,}01}{0{,}9\cdot0{,}01+0{,}1\cdot0{,}99}
        =\frac{0{,}009}{0{,}108}=\frac1{12}\approx0{,}083,
\]
\[
  \eta(0)=\frac{0{,}1\cdot0{,}01}{0{,}1\cdot0{,}01+0{,}9\cdot0{,}99}
        =\frac{0{,}001}{0{,}892}=\frac1{892}\approx0{,}0011 .
\]
Um resultado positivo multiplica a probabilidade de doença por cerca de $8$, de $1\%$ para
$8{,}3\%$, e um negativo a divide por cerca de $9$; ainda assim, mesmo com o exame positivo,
a doença continua improvável, porque ela é rara: entre os positivos, os $9$ doentes em cada
mil pessoas são superados pelos $99$ saudáveis com alarme falso.

\smallskip
\textbf{(b)} O classificador de Bayes declara doente quando $\eta(x)>1/2$. Como
$\eta(0)$ e $\eta(1)$ estão abaixo de $1/2$, ele é $g\equiv0$: declara todos saudáveis,
qualquer que seja o resultado do exame. A probabilidade de erro dele é
$\Prob(Y=1)=0{,}01$. Seguir o exame erra nos doentes com $X=0$ e nos saudáveis com $X=1$:
\[
  \Prob(X\ne Y)=0{,}01\cdot0{,}1+0{,}99\cdot0{,}1=0{,}001+0{,}099=0{,}1 ,
\]
dez vezes mais. A conclusão do colega confunde a probabilidade de erro com a utilidade. O
classificador de Bayes minimiza a probabilidade de erro, que dá o mesmo peso a um doente não
detectado e a um alarme falso; com uma doença rara, a maneira mais barata de errar pouco é
nunca acusar ninguém, ao preço de não detectar doente nenhum. O exame detecta $90\%$ dos
doentes: o que falta saber é quanto vale cada um deles comparado a um alarme falso.

\smallskip
\textbf{(c)} O corte é $\ell_1/(\ell_0+\ell_1)=1/(1+r)$, e o classificador compara $\eta(x)$
com ele em cada valor de $x$ separadamente. Ele segue o exame, isto é, declara doente em
$x=1$ e saudável em $x=0$, quando
\[
  \eta(1)>\frac1{1+r}\iff1+r>12\iff r>11
  \qquad\text{e}\qquad
  \eta(0)\le\frac1{1+r}\iff r\le891 .
\]
Com $r\le11$, ele declara todos saudáveis; com $r>891$, declara todos doentes, porque,
quando um doente não detectado custa mais que $891$ alarmes falsos, até a chance de $1/892$
que resta depois de um exame negativo justifica declarar doente. Nos extremos da faixa há
empate: com $r=11$, declarar doente ou saudável custa o mesmo em $x=1$, e, com $r=891$, em
$x=0$. Nos dois empates o classificador declara saudável, porque só declara doente quando
$\eta(x)$ passa do corte; por isso a faixa exclui $r=11$ e inclui $r=891$. Em termos de
custo, seguir o exame é uma decisão ótima para $11\le r\le891$, e a única para
$11<r<891$: vale a pena exatamente quando deixar passar um doente custa entre $11$ e $891$
alarmes falsos.

Com $\ell_0=20$ e $\ell_1=1$, $r=20$ está na faixa, e seguir o exame é o ótimo. Numa
população de $10\,000$ pessoas há, em média, $P=100$ doentes e $N=9\,900$ saudáveis, e as
contagens esperadas saem das taxas: $\mathrm{FN}=(1-\mathrm{TPR})\,P$ e
$\mathrm{FP}=\mathrm{FPR}\cdot N$. Declarar todos saudáveis dá $\mathrm{FN}=100$ e custo
$20\cdot100=2\,000$. Seguir o exame dá $\mathrm{FN}=0{,}1\cdot100=10$ e
$\mathrm{FP}=0{,}1\cdot9\,900=990$, com custo $20\cdot10+990=1\,190$. Declarar todos doentes dá
$\mathrm{FP}=9\,900$ e custo $9\,900$. O exame economiza $810$ unidades de custo a cada dez mil
pessoas em relação ao classificador de Bayes, embora erre dez vezes mais que ele.
\end{solucao}

\end{exercicio}

%==============================================================================
\section*{Aula 09 --- Máquinas de Vetores de Suporte}
%==============================================================================

\begin{exercicio}
\label{ex:09-suporte}
% Aula 09 --- a margem máxima depende só dos vetores de suporte, e o LOOCV fica limitado por |S|/n.
Tiago trabalha no controle de qualidade de uma fábrica e classifica peças em boas ($y=+1$) e
defeituosas ($y=-1$) a partir de $p$ medidas. As suas $n$ observações de treino
$(\x_i,y_i)$ são linearmente separáveis, e cada classe tem pelo menos duas observações. O
classificador de margem máxima é dado pela solução do problema
\[
  \max_{\beta_0,\,\bbeta,\,M}\ M
  \quad\text{sujeito a}\quad
  \sum_{j=1}^p\beta_j^2=1
  \quad\text{e}\quad
  y_i\,(\beta_0+\bbeta^\top\x_i)\ge M,\ \ i=1,\dots,n .
\]
Escreva $f(\x)=\beta_0+\bbeta^\top\x$. Os vetores de suporte são as observações com
$y_if(\x_i)=M$, e $\mathcal S$ é o conjunto dos seus índices. Admita que, para observações
separáveis com as duas classes presentes, o problema tem solução, e seja $M^*>0$ o valor
ótimo. As notas afirmam, sem demonstrar, que essa solução depende apenas dos vetores de
suporte; Tiago quer a demonstração e quer saber o que ela lhe dá.
\begin{itens}
  \item \pts{1,0} Mostre que a solução é única. (Se houvesse duas soluções com vetores
        $\bbeta$ diferentes, compare-as com o ponto médio entre elas. Se tivessem o mesmo
        $\bbeta$, olhe, como funções de $\beta_0$, o menor valor de $y_if(\x_i)$ em cada
        classe.)
  \item \pts{1,4} Seja $i_0\notin\mathcal S$. Mostre que o classificador de margem máxima
        das $n-1$ observações sem a observação $i_0$ é o mesmo das $n$ observações. (Se o
        valor ótimo sem $i_0$ fosse maior que $M^*$, combine a solução nova com a antiga, com
        peso pequeno na nova.) Mostre também que deslocar $\x_{i_0}$ para qualquer ponto
        $\x'$ com $y_{i_0}f(\x')>M^*$, em que $f$ é a função da solução, não muda o
        classificador.
  \item \pts{1,0} A estimativa do risco por validação cruzada com uma observação deixada de
        fora (LOOCV) é a proporção das observações $i$ que o classificador ajustado sem $i$
        classifica errado. Mostre que, para o classificador de margem máxima, ela é no
        máximo $\abs{\mathcal S}/n$.
  \item \pts{0,6} Tiago ajusta o classificador a $60$ peças e encontra $3$ vetores de
        suporte. O que ele pode afirmar sobre a estimativa LOOCV sem fazer mais nenhum ajuste?
        Ele refaz os $60$ ajustes e não encontra erro nenhum. Isso contradiz o item (c)? O que
        a cota não garante?
\end{itens}

\begin{solucao}
\textbf{(a)} Sejam $(\beta_0,\bbeta,M^*)$ e $(\beta_0',\bbeta',M^*)$ duas soluções, e
$(\bar b,\bar{\bm w})$ o ponto médio de $(\beta_0,\bbeta)$ e $(\beta_0',\bbeta')$. Para cada
$i$,
\[
  y_i\big(\bar b+\bar{\bm w}^\top\x_i\big)
  =\tfrac12\big[y_i(\beta_0+\bbeta^\top\x_i)+y_i(\beta_0'+\bbeta'^\top\x_i)\big]\ \ge\ M^* .
\]
Suponha $\bbeta\ne\bbeta'$. Como $\norm{\bbeta}=\norm{\bbeta'}=1$, a desigualdade de
Cauchy--Schwarz dá $\bbeta^\top\bbeta'<1$, com desigualdade estrita porque a igualdade
exigiria $\bbeta'=\bbeta$, e
$\norm{\bar{\bm w}}^2=\tfrac14\big(2+2\,\bbeta^\top\bbeta'\big)<1$. Também
$\bar{\bm w}\ne\bm0$: do contrário, $y_i\bar b\ge M^*>0$ para todo $i$, o que exigiria
$\bar b>0$ nas observações da classe $+1$ e $\bar b<0$ nas da classe $-1$. Dividindo
$(\bar b,\bar{\bm w})$ por $\norm{\bar{\bm w}}\in(0,1)$, obtém-se um vetor de norma $1$ com
$y_i(\cdot)\ge M^*/\norm{\bar{\bm w}}>M^*$ para todo $i$, o que contradiz a otimalidade de
$M^*$. Logo $\bbeta=\bbeta'$.

Com $\bbeta$ fixado, sejam
$m_+(b)=\min_{i:\,y_i=+1}(b+\bbeta^\top\x_i)$ e $m_-(b)=\min_{i:\,y_i=-1}(-b-\bbeta^\top\x_i)$,
de modo que a margem do hiperplano $(b,\bbeta)$ é $\min\{m_+(b),m_-(b)\}$. A função $m_+$ é
estritamente crescente, e $m_-$, estritamente decrescente. Suponha $\beta_0<\beta_0'$. Então
$m_-(\beta_0)>m_-(\beta_0')\ge M^*$, e, como o mínimo em $\beta_0$ vale $M^*$, é
$m_+(\beta_0)=M^*$; do mesmo modo, $m_+(\beta_0')>m_+(\beta_0)=M^*$ e $m_-(\beta_0')=M^*$. No
ponto médio $b=(\beta_0+\beta_0')/2$, $m_+(b)>m_+(\beta_0)=M^*$ e $m_-(b)>m_-(\beta_0')=M^*$:
o hiperplano $(b,\bbeta)$ teria margem maior que $M^*$, um absurdo. Logo $\beta_0=\beta_0'$, e
a solução é única.

\smallskip
\textbf{(b)} Seja $(\beta_0^*,\bbeta^*,M^*)$ a solução com as $n$ observações e $f^*$ a função
correspondente; como $i_0\notin\mathcal S$, vale $\delta=y_{i_0}f^*(\x_{i_0})-M^*>0$. Sem
$i_0$, as observações continuam separáveis e com as duas classes, porque cada classe tinha
pelo menos duas; o problema tem solução, com valor ótimo $M'$. O ponto
$(\beta_0^*,\bbeta^*,M^*)$ satisfaz todas as restrições sem $i_0$, e portanto $M'\ge M^*$.

Suponha $M'>M^*$, atingido por $(\beta_0',\bbeta')$, com função $f'$. Para $t\in(0,1]$, seja
$(b_t,\bm w_t)=(1-t)(\beta_0^*,\bbeta^*)+t(\beta_0',\bbeta')$. Para $i\ne i_0$,
\[
  y_i\big(b_t+\bm w_t^\top\x_i\big)=(1-t)\,y_if^*(\x_i)+t\,y_if'(\x_i)\ \ge\ (1-t)M^*+tM'\ >\ M^* ,
\]
e, para $i_0$, com $c=\abs{y_{i_0}f'(\x_{i_0})-y_{i_0}f^*(\x_{i_0})}$,
\[
  y_{i_0}\big(b_t+\bm w_t^\top\x_{i_0}\big)\ \ge\ M^*+\delta-t\,c\ >\ M^*
  \qquad\text{para }0<t<\delta/c
\]
(para todo $t$, se $c=0$). Pela desigualdade triangular,
$\norm{\bm w_t}\le(1-t)\norm{\bbeta^*}+t\norm{\bbeta'}=1$, e $\bm w_t\ne\bm0$ pelo mesmo
argumento do item (a). Fixe um tal $t$ e divida $(b_t,\bm w_t)$ por $\norm{\bm w_t}\le1$: os
números $y_i(\cdot)$, todos maiores que $M^*>0$, não diminuem, e o resultado é um hiperplano
de norma $1$ com $y_if(\x_i)>M^*$ para \emph{todas} as $n$ observações. Isso contradiz a
otimalidade de $M^*$. Logo $M'=M^*$, e $(\beta_0^*,\bbeta^*,M^*)$, que é viável sem $i_0$ e
atinge esse valor, é uma solução sem $i_0$; pelo item (a), aplicado às $n-1$ observações, é a
única.

Para o deslocamento, considere as $n-1$ observações sem $i_0$ mais a observação
$(\x',y_{i_0})$. Todo ponto viável desse problema é viável sem $i_0$, e por isso o valor ótimo
é no máximo $M^*$. A solução $(\beta_0^*,\bbeta^*,M^*)$ satisfaz as restrições sem $i_0$ e a
nova, $y_{i_0}f^*(\x')>M^*$: é viável e atinge $M^*$, logo é ótima e, pelo item (a), a única. A
solução não vê as observações que estão do lado certo e fora da faixa: elas podem sair ou
mudar de lugar, desde que continuem fora dela.

\smallskip
\textbf{(c)} Para cada $i$, as observações sem $i$ são separáveis e têm as duas classes, e o
classificador de margem máxima delas, $g_{-i}$, está bem definido. Se $i\notin\mathcal S$,
pelo item (b) $g_{-i}$ é o classificador $g$ das $n$ observações, que classifica $\x_i$
corretamente, porque $y_if^*(\x_i)>M^*>0$. Assim, $g_{-i}(\x_i)\ne y_i$ só pode acontecer para
$i\in\mathcal S$, e
\[
  \frac1n\sum_{i=1}^n\1{g_{-i}(\x_i)\ne y_i}
  =\frac1n\sum_{i\in\mathcal S}\1{g_{-i}(\x_i)\ne y_i}\ \le\ \frac{\abs{\mathcal S}}{n} .
\]

\smallskip
\textbf{(d)} Com um único ajuste, Tiago sabe que a estimativa LOOCV é no máximo
$3/60=0{,}05$: das $60$ retiradas, $57$ deixam o classificador intacto e não erram, e só as
$3$ de vetores de suporte podem errar. Os $60$ ajustes, que mostram erro zero, não contradizem
nada: a cota é uma desigualdade, e os vetores de suporte deixados de fora podem continuar do
lado certo do hiperplano recalculado. Neste conjunto de dados, nenhum deles muda de lado.

A cota limita uma estimativa do risco, não o risco. Ela não diz que peças futuras serão
classificadas sem erro: o treino ser separável não significa que a população seja, e uma peça
nova perto da fronteira pode deslocar muito o hiperplano, que é a sensibilidade a observações
novas apontada nas notas. O que ela mostra é que a estimativa LOOCV deste classificador
custa um ajuste em vez de $n$, desde que se aceite uma cota no lugar do número.
\end{solucao}

\end{exercicio}

\begin{exercicio}
\label{ex:09-radial}
% Aula 09 --- produtos e exponenciais de kernels; o kernel radial em qualquer dimensão.
Júlia estuda matemática aplicada e lê nas notas que o \emph{kernel} radial
$k(\x,\x')=\exp\!\big(-\gamma\norm{\x-\x'}^2\big)$, com $\gamma>0$, é um \emph{kernel} em
$\R^p$; o exemplo das notas, porém, só exibe a função $\varphi$ para $p=1$. Ela decide provar
o caso geral. Lembre que $k:\R^p\times\R^p\to\R$ é um \emph{kernel} se existem um espaço de
Hilbert $\mathcal H$ e uma função $\varphi:\R^p\to\mathcal H$ com
$k(\x,\x')=\inner{\varphi(\x)}{\varphi(\x')}_{\mathcal H}$ para todos $\x,\x'$, e que, pelo
teorema visto em aula, isso acontece se e somente se $k$ é simétrica e, para todo $n\ge1$ e
quaisquer $\x_1,\dots,\x_n\in\R^p$, a matriz $\big[k(\x_i,\x_j)\big]_{i,j=1}^n$ é positiva
semidefinida. Use à vontade que o produto interno $\inner{\x}{\x'}$ e as funções constantes
iguais a $c\ge0$ são \emph{kernels}, e que a soma de dois \emph{kernels} e o produto de um
\emph{kernel} por uma constante $c\ge0$ também são.
\begin{itens}
  \item \pts{1,2} Mostre que, se $k_1$ e $k_2$ são \emph{kernels}, o produto
        $k(\x,\x')=k_1(\x,\x')\,k_2(\x,\x')$ também é. (Escreva uma das duas matrizes pela
        decomposição espectral.)
  \item \pts{1,0} Mostre que, se $k$ é um \emph{kernel} e $h:\R^p\to\R$ é uma função
        qualquer, então $h(\x)\,h(\x')\,k(\x,\x')$ é um \emph{kernel}. Mostre também que
        $e^{k(\x,\x')}$ é um \emph{kernel}, usando a série da exponencial e o fato de que um
        limite de matrizes positivas semidefinidas é positivo semidefinido.
  \item \pts{0,8} Conclua que o \emph{kernel} radial é um \emph{kernel} em $\R^p$, para todo
        $p$ e todo $\gamma>0$. (Escreva $\norm{\x-\x'}^2$ em termos de $\norm{\x}^2$,
        $\norm{\x'}^2$ e $\inner{\x}{\x'}$.)
  \item \pts{1,0} Mostre que, para todo \emph{kernel},
        $\norm{\varphi(\x)-\varphi(\x')}_{\mathcal H}^2=k(\x,\x)-2\,k(\x,\x')+k(\x',\x')$, e
        calcule essa distância para o \emph{kernel} radial. Conclua que $\varphi$ leva $\R^p$
        na esfera unitária de $\mathcal H$ e que pontos distintos ficam, em $\mathcal H$, a
        distância menor que $\sqrt2$. Com $\gamma=0{,}5$ e com $\gamma=50$, a que distância
        dois pontos precisam estar em $\R^p$ para ficarem à distância $1$ em $\mathcal H$?
        Interprete o resultado em termos da flexibilidade da SVM. \emph{Dado:}
        $\sqrt{2\log2}\approx1{,}177$, com $\log$ o logaritmo natural.
\end{itens}

\begin{solucao}
\textbf{(a)} Sejam $\x_1,\dots,\x_n\in\R^p$, $A=\big[k_1(\x_i,\x_j)\big]$ e
$B=\big[k_2(\x_i,\x_j)\big]$, simétricas e positivas semidefinidas pelo teorema. A matriz de
$k$ nesses pontos é o produto entrada a entrada, $(A\circ B)_{ij}=A_{ij}B_{ij}$, que é
simétrica. Pelo teorema espectral, $B=\sum_{m=1}^n\lambda_m\bm v_m\bm v_m^\top$, com
autovalores $\lambda_m\ge0$, porque $B$ é positiva semidefinida, e autovetores ortonormais
$\bm v_m\in\R^n$; em coordenadas, $B_{ij}=\sum_m\lambda_m v_{mi}v_{mj}$. Para $\bm a\in\R^n$,
\begin{align*}
  \bm a^\top(A\circ B)\,\bm a
  &=\sum_{i,j}a_ia_jA_{ij}\sum_m\lambda_mv_{mi}v_{mj}
   =\sum_m\lambda_m\sum_{i,j}(a_iv_{mi})\,A_{ij}\,(a_jv_{mj})\\
  &=\sum_m\lambda_m\,\bm u_m^\top A\,\bm u_m\ \ge\ 0,
\end{align*}
em que $\bm u_m=(a_1v_{m1},\dots,a_nv_{mn})$, e cada parcela é não negativa porque
$\lambda_m\ge0$ e $A$ é positiva semidefinida. Como isso vale para todo $n$ e quaisquer
pontos, o teorema dá que $k_1k_2$ é um \emph{kernel}.

\smallskip
\textbf{(b)} A função $h(\x)h(\x')k(\x,\x')$ é simétrica, e a sua matriz nos pontos
$\x_1,\dots,\x_n$ é $DKD$, com $K=\big[k(\x_i,\x_j)\big]$ e
$D=\operatorname{diag}\big(h(\x_1),\dots,h(\x_n)\big)$. Para $\bm a\in\R^n$,
$\bm a^\top DKD\,\bm a=(D\bm a)^\top K(D\bm a)\ge0$. (Pela definição, é o mesmo que tomar a
função $h\varphi$ no lugar de $\varphi$.)

Para a exponencial, $e^{k(\x,\x')}=\sum_{m\ge0}k(\x,\x')^m/m!$. A potência $k^0$ é a
constante $1$, que é um \emph{kernel}, e, se $k^m$ é um \emph{kernel}, $k^{m+1}=k^m\cdot k$
também é, pelo item (a); por indução, todas as potências são \emph{kernels}. As somas
parciais $S_M=\sum_{m=0}^Mk^m/m!$ são, então, \emph{kernels}, como somas de \emph{kernels}
multiplicados por constantes positivas. Fixados os pontos e $\bm a\in\R^n$, cada número
$\bm a^\top\big[S_M(\x_i,\x_j)\big]\bm a$ é não negativo, e, quando $M\to\infty$, converge para
$\bm a^\top\big[e^{k(\x_i,\x_j)}\big]\bm a$, porque é uma soma finita de termos que convergem.
Um limite de números não negativos é não negativo. A matriz de $e^k$ é, portanto, positiva
semidefinida, e $e^k$, que é simétrica, é um \emph{kernel}.

\smallskip
\textbf{(c)} Como $\norm{\x-\x'}^2=\norm{\x}^2-2\inner{\x}{\x'}+\norm{\x'}^2$,
\[
  e^{-\gamma\norm{\x-\x'}^2}
  =e^{-\gamma\norm{\x}^2}\,e^{-\gamma\norm{\x'}^2}\,e^{2\gamma\inner{\x}{\x'}} .
\]
O produto interno é um \emph{kernel}, e $2\gamma\inner{\x}{\x'}$ também, porque $2\gamma>0$.
Pelo item (b), $e^{2\gamma\inner{\x}{\x'}}$ é um \emph{kernel}, e, ainda pelo item (b), com
$h(\x)=e^{-\gamma\norm{\x}^2}$, o produto acima é um \emph{kernel}. O argumento vale para todo
$p$ e todo $\gamma>0$, e não precisa exibir $\varphi$; com $p=1$, expandir
$e^{2\gamma xx'}$ em série reproduz a função $\varphi$ do exemplo das notas.

\smallskip
\textbf{(d)} Pela bilinearidade e pela simetria do produto interno de $\mathcal H$,
\begin{align*}
  \norm{\varphi(\x)-\varphi(\x')}^2
  &=\inner{\varphi(\x)}{\varphi(\x)}-2\inner{\varphi(\x)}{\varphi(\x')}+\inner{\varphi(\x')}{\varphi(\x')}\\
  &=k(\x,\x)-2\,k(\x,\x')+k(\x',\x') .
\end{align*}
No \emph{kernel} radial, $k(\x,\x)=e^0=1$, de modo que $\norm{\varphi(\x)}=1$ para todo $\x$,
e
\[
  \norm{\varphi(\x)-\varphi(\x')}^2=2\big(1-e^{-\gamma\norm{\x-\x'}^2}\big) .
\]
Essa quantidade é nula se e somente se $\x=\x'$, cresce com $\norm{\x-\x'}$ e é sempre menor
que $2$: os pontos vão para a esfera unitária de $\mathcal H$, a distâncias menores que
$\sqrt2$, isto é, com produto interno $k(\x,\x')>0$ e ângulo menor que $90^\circ$.

Dois pontos ficam à distância $1$ em $\mathcal H$ quando
$2\big(1-e^{-\gamma\norm{\x-\x'}^2}\big)=1$, isto é, quando $k(\x,\x')=1/2$, a meio caminho
entre pontos iguais ($k=1$) e ortogonais ($k=0$). Isso acontece se e somente se
\[
  \norm{\x-\x'}=\sqrt{\frac{\log2}{\gamma}} ,
\]
que vale $\sqrt{2\log2}\approx1{,}177$ com $\gamma=0{,}5$ e $\sqrt{2\log2}/10\approx0{,}118$
com $\gamma=50$, dez vezes menos.

A distância em $\mathcal H$ depende de $\x$ e $\x'$ só por meio de
$\sqrt\gamma\,\norm{\x-\x'}$: multiplicar $\gamma$ por $100$ tem o mesmo efeito que
multiplicar todas as covariáveis por $10$, e $1/\sqrt\gamma$ faz o papel de uma unidade de
comprimento em $\R^p$. É também a escala em que age cada parcela da função de decisão
$f(\x)=\beta_0+\sum_{i\in\mathcal S}\alpha_ik(\x,\x_i)$: $k(\x,\x_i)$ cai à metade quando
$\norm{\x-\x_i}=\sqrt{\log2/\gamma}$ e decai depressa a partir daí. Com $\gamma=50$, a
influência de cada vetor de suporte se restringe a uma vizinhança de raio da ordem de
$0{,}1$ em torno dele, a fronteira pode contornar observações isoladas, e o ajuste fica
flexível, com risco de superajuste. Com $\gamma=0{,}5$, o raio é dez vezes maior, muitos
vetores de suporte pesam em cada ponto, e a fronteira é suave.
\end{solucao}

\end{exercicio}

%==============================================================================
\section*{Aula 10 --- KNN e Árvores de Classificação}
%==============================================================================

\begin{exercicio}
\label{ex:10-voto}
% Aula 10 --- o voto de k vizinhos: a variância cai com k, e o 1-NN erra no máximo o dobro do classificador de Bayes.
Otávio é agrônomo e usa o KNN para classificar os talhões de uma fazenda em atacados por uma
praga ($Y=1$) ou não ($Y=0$), a partir de medidas de solo e de clima. Num talhão novo, com
covariáveis $\x$, a probabilidade de praga é $\eta(\x)=\Prob(Y=1\mid\X=\x)=0{,}6$. O KNN
estima essa probabilidade pela proporção de praga entre os $k$ talhões de treino mais
próximos de $\x$, cujos índices formam o conjunto $\mathcal{N}_{\x}^{(k)}$,
\[
  \widehat P=\frac1k\sum_{i\in\mathcal{N}_{\x}^{(k)}}Y_i ,
\]
e declara praga quando $\widehat P>1/2$. Tome $k$ ímpar, para que não haja empate. Para
isolar o papel de $k$, suponha que os rótulos dos vizinhos sejam independentes entre si e
do rótulo $Y$ do talhão novo, e que cada um seja igual a $1$ com probabilidade $0{,}6$, como
se todos os vizinhos estivessem no próprio $\x$.
\begin{itens}
  \item \pts{1,1} Mostre que $\E[\widehat P]=0{,}6$ e $\Var(\widehat P)=0{,}24/k$.
  \item \pts{1,8} Seja $q_k=\Prob(\widehat P>1/2)$. Mostre que a probabilidade de o KNN errar
        o talhão novo é $0{,}6-0{,}2\,q_k$. Calcule $q_5$ e, usando $q_1=0{,}600$,
        $q_{15}\approx0{,}787$ e $q_{51}\approx0{,}926$, a probabilidade de erro para
        $k=1,5,15$ e $51$. Use a desigualdade de Chebyshev para mostrar que $q_k\to1$ e que a
        probabilidade de erro tende a $0{,}4$. Que classificador erra com probabilidade
        $0{,}4$ nesse talhão?
  \item \pts{1,1} Mostre que, com $k=1$ e uma probabilidade $\eta$ qualquer no lugar de
        $0{,}6$, a probabilidade de erro é $2\eta(1-\eta)$, e que
        $\min(\eta,1-\eta)\le2\eta(1-\eta)\le2\min(\eta,1-\eta)$. Interprete. O que acontece
        no outro extremo, com $k=n$?
\end{itens}

\begin{solucao}
\textbf{(a)} O conjunto $\mathcal{N}_{\x}^{(k)}$ tem $k$ elementos, e cada $Y_i$, com
$i\in\mathcal{N}_{\x}^{(k)}$, é uma Bernoulli$(0{,}6)$, com média $0{,}6$ e variância
$0{,}6\cdot0{,}4=0{,}24$. A esperança sai da linearidade, e a variância, da independência dos
rótulos, que permite somar as variâncias:
\[
  \E[\widehat P]=\frac1k\sum_{i\in\mathcal{N}_{\x}^{(k)}}\E[Y_i]=\frac{k\cdot0{,}6}{k}=0{,}6,
  \qquad
  \Var(\widehat P)=\frac1{k^2}\sum_{i\in\mathcal{N}_{\x}^{(k)}}\Var(Y_i)
  =\frac{k\cdot0{,}24}{k^2}=\frac{0{,}24}{k} .
\]
Nesta idealização, a proporção de vizinhos é um estimador não viesado de $\eta(\x)$, e a sua
variância cai como $1/k$.

\smallskip
\textbf{(b)} O KNN erra se declara praga num talhão sem praga ou deixa de declarar num talhão
com praga. Como a decisão depende só dos rótulos dos vizinhos, que são independentes de $Y$,
\[
  \Prob(\ghat(\x)\ne Y)=q_k\,\Prob(Y=0)+(1-q_k)\,\Prob(Y=1)=0{,}4\,q_k+0{,}6\,(1-q_k)=0{,}6-0{,}2\,q_k .
\]
Com $k=5$, o KNN declara praga quando pelo menos $3$ dos $5$ vizinhos têm praga, e o número
deles é uma binomial$(5;\,0{,}6)$:
\[
  q_5=\binom53(0{,}6)^3(0{,}4)^2+\binom54(0{,}6)^4(0{,}4)+(0{,}6)^5
     =0{,}3456+0{,}2592+0{,}0778\approx0{,}683 .
\]
As probabilidades de erro são $0{,}480$, $0{,}463$, $0{,}443$ e $0{,}415$ para $k=1$, $5$, $15$ e
$51$. Como $k$ é ímpar, $\widehat P\ne1/2$, e o KNN deixa de declarar praga exatamente
quando $\widehat P<1/2$. Esse evento está contido em
$\{\abs{\widehat P-0{,}6}\ge0{,}1\}$, e, pela desigualdade de Chebyshev,
\[
  1-q_k=\Prob\big(\widehat P<1/2\big)\le\Prob\big(\abs{\widehat P-0{,}6}\ge0{,}1\big)
  \le\frac{\Var(\widehat P)}{0{,}1^2}=\frac{24}{k}\ \longrightarrow\ 0 .
\]
Logo $q_k\to1$ e a probabilidade de erro tende a $0{,}6-0{,}2=0{,}4$. É a probabilidade de
erro do classificador de Bayes nesse talhão: como $\eta(\x)>1/2$, ele declara praga, e erra
com probabilidade $1-\eta(\x)=0{,}4$, que nenhum classificador pode melhorar.

\smallskip
\textbf{(c)} Com $k=1$, o KNN repete o rótulo do vizinho mais próximo, que é $1$ com
probabilidade $\eta$ e independente de $Y$. Ele erra se os dois rótulos diferem:
\[
  \Prob(\ghat(\x)\ne Y)=\eta(1-\eta)+(1-\eta)\eta=2\eta(1-\eta) .
\]
Sejam $m=\min(\eta,1-\eta)$ e $1-m=\max(\eta,1-\eta)$, com $1/2\le1-m\le1$. Então
$2\eta(1-\eta)=2m(1-m)$, e multiplicar as desigualdades $1/2\le1-m\le1$ por $2m\ge0$ dá
$m\le2m(1-m)\le2m$. Como $m$ é a probabilidade de erro do classificador de Bayes em $\x$, o
$1$-NN, nesta idealização, erra pelo menos tanto quanto ele e no máximo o dobro; com
$\eta=0{,}6$, são $0{,}48$ contra $0{,}40$. No outro extremo, com $k=n$, todos os pontos de treino
votam, $\widehat P$ é a proporção de talhões com praga no treino inteiro, e o KNN atribui a
mesma classe, a majoritária, a qualquer talhão: é o classificador constante, com o maior viés
possível.
\end{solucao}

\end{exercicio}

\begin{exercicio}
\label{ex:10-poda}
% Aula 10 --- crescer com Gini, podar com o erro: a poda remove os ramos que não mudam nenhuma decisão.
Rafael é analista numa operadora de planos de saúde e cresceu uma árvore de classificação,
com o índice de Gini, para prever quais beneficiários serão internados no próximo ano
($Y=1$). Agora vai podá-la. Para uma região $R$ com $n_R$ observações de treino, sejam
$n_{R,c}$ o número delas na classe $c$ e $\widehat p_{R,c}=n_{R,c}/n_R$; o índice de Gini de
$R$ é $G(R)=\sum_c\widehat p_{R,c}(1-\widehat p_{R,c})$. A folha $R$ prediz a classe
majoritária, e $E(R)=1-\max_c\widehat p_{R,c}$ é o seu erro de classificação, de modo
que $n_RE(R)=n_R-\max_cn_{R,c}$ é o número de erros de treino na folha. Na poda por
custo-complexidade com o erro de classificação, escolhe-se, entre as subárvores $T$ da árvore
crescida, a que minimiza
\[
  C_\alpha(T)=\sum_{m=1}^{\abs{T}}n_mE(R_m)+\alpha\abs{T},
\]
em que $R_1,\dots,R_{\abs{T}}$ são as folhas de $T$, $n_m=n_{R_m}$ e $\abs{T}$ é o número de
folhas; o valor de $\alpha>0$ é escolhido por validação cruzada.
\begin{itens}
  \item \pts{1,3} Seja $R$ uma região dividida em regiões disjuntas $R_1,\dots,R_L$. Mostre
        que $\sum_ln_{R_l}E(R_l)\le n_RE(R)$, com qualquer número de classes, e que vale a
        igualdade quando uma mesma classe $c^*$ é majoritária em todas as $R_l$; nesse caso,
        $c^*$ é majoritária também em $R$.
  \item \pts{1,6} Mostre que, para todo $\alpha>0$, nenhuma subárvore que minimize $C_\alpha$
        tem um nó interno cujas folhas descendentes predizem todas a mesma classe.
  \item \pts{1,1} Na árvore de Rafael, um nó $R$ com $400$ beneficiários, dos quais $100$
        foram internados, foi dividido em dois filhos de $200$, que são folhas da árvore
        crescida: $R_1$, com $20$ internados, e $R_2$, com $80$. Calcule a redução do índice
        de Gini ponderado,
        $G(R)-\big[n_1G(R_1)+n_2G(R_2)\big]/n_R$, e a redução do erro de classificação
        ponderado. Por que o crescimento fez essa divisão, e o que a poda fará com ela?
\end{itens}

\begin{solucao}
\textbf{(a)} As regiões $R_l$ particionam $R$, e por isso $n_{R,c}=\sum_ln_{R_l,c}$ para
cada classe $c$. O máximo de uma soma não passa da soma dos máximos:
\[
  \max_cn_{R,c}=\max_c\sum_ln_{R_l,c}\ \le\ \sum_l\max_cn_{R_l,c} .
\]
Logo
\[
  \sum_ln_{R_l}E(R_l)=\sum_l\Big(n_{R_l}-\max_cn_{R_l,c}\Big)=n_R-\sum_l\max_cn_{R_l,c}
  \ \le\ n_R-\max_cn_{R,c}=n_RE(R).
\]
Se $c^*$ é majoritária em todas as $R_l$, então
$\sum_l\max_cn_{R_l,c}=\sum_ln_{R_l,c^*}=n_{R,c^*}\le\max_cn_{R,c}$; com a desigualdade
anterior, os três números são iguais. A igualdade vale, e $n_{R,c^*}=\max_cn_{R,c}$: $c^*$ é
majoritária em $R$. Dividir uma região nunca aumenta o número de erros de treino, e só o
diminui quando alguma das partes muda a classe predita.

\smallskip
\textbf{(b)} Suponha que $T$ minimize $C_\alpha$ e tenha um nó interno $t$ cujas
$L\ge2$ folhas descendentes predizem todas a classe $c^*$. Seja $T'$ a subárvore obtida
trocando o ramo que sai de $t$ por uma única folha, a região de $t$. Pelo item (a), essa folha
tem $c^*$ como classe majoritária e o mesmo número de erros de treino que as $L$ folhas
juntas; as outras folhas não mudam. As duas subárvores cometem, então, os mesmos erros, e
$\abs{T'}=\abs{T}-(L-1)$, de modo que
\[
  C_\alpha(T')=C_\alpha(T)-\alpha(L-1)<C_\alpha(T),
\]
porque $\alpha>0$ e $L\ge2$. Isso contradiz a escolha de $T$. Qualquer que seja $\alpha>0$, a
poda pelo erro de classificação remove todos os ramos que não mudam nenhuma decisão.

\smallskip
\textbf{(c)} As proporções de internados são $0{,}25$ no nó e $0{,}1$ e $0{,}4$ nos filhos. Com
duas classes, as duas parcelas de $G$ são iguais, e $G=2\widehat p(1-\widehat p)$, com
$\widehat p$ a proporção de internados:
\[
  G(R)=2\cdot0{,}25\cdot0{,}75=0{,}375,
  \qquad
  \tfrac12\,G(R_1)+\tfrac12\,G(R_2)=\tfrac12\cdot0{,}18+\tfrac12\cdot0{,}48=0{,}33 ,
\]
uma redução de $0{,}045$. A classe majoritária é ``não internado'' no nó ($300$ contra $100$) e
nos dois filhos ($180$ contra $20$ e $120$ contra $80$). Os erros são $100$ no nó e
$20+80=100$ nos filhos: a redução do erro é zero, como o item (a) previa. O Gini é
estritamente côncavo e registra o ganho de pureza de separar um grupo de risco baixo
($10\%$ de internações) de um de risco alto ($40\%$); por isso o crescimento fez a divisão.
Como os filhos são folhas da árvore crescida, numa subárvore que mantenha a divisão eles são
as únicas folhas descendentes de $R$, e os dois predizem a mesma classe. Pelo item (b), a poda
pelo erro remove a divisão, qualquer que seja $\alpha>0$.
\end{solucao}

\end{exercicio}

\begin{exercicio}
\label{ex:10-impureza}
% Aula 10 --- uma divisão em que o Gini e a entropia discordam, e a forma da curva que explica isso.
Ana é bióloga e usa uma árvore de classificação para distinguir mudas de duas espécies, a
espécie $1$ e a espécie $0$, a partir de características medidas no viveiro. Num nó com $12$
mudas, $6$ de cada espécie, a árvore compara duas divisões:
\begin{description}[leftmargin=1.4cm,style=nextline]
  \item[A] pela altura, separa dois grupos de $6$ mudas, um com $2$ da espécie $1$ e $4$ da
        espécie $0$ e outro com $4$ da espécie $1$ e $2$ da espécie $0$;
  \item[B] por uma mancha rara na folha, isola uma única muda, da espécie $0$, das outras $11$,
        que têm $6$ da espécie $1$ e $5$ da espécie $0$.
\end{description}
Numa região com proporção $p$ de mudas da espécie $1$, o índice de Gini é $G(p)=2p(1-p)$, a
entropia é $D(p)=-p\log_2p-(1-p)\log_2(1-p)$, com $0\log_20=0$, e o erro de classificação é
$E(p)=\min(p,1-p)$. Uma divisão em regiões com $n_1$ e $n_2$ mudas e proporções $p_1$ e $p_2$
tem impureza ponderada $\big[n_1I(p_1)+n_2I(p_2)\big]/(n_1+n_2)$, para cada medida $I$, e a
árvore escolhe a divisão de menor impureza ponderada.
\begin{itens}
  \item \pts{2,0} Calcule a impureza ponderada das divisões A e B pelas três medidas e diga
        qual divisão cada medida escolhe. \emph{Dados:} $\log_23\approx1{,}5850$,
        $\log_25\approx2{,}3219$ e $\log_211\approx3{,}4594$.
  \item \pts{2,0} Seja $I$ uma medida de impureza qualquer para duas classes, com
        $I(p)=I(1-p)$, $I(0)=0$ e $I(6/11)>0$. Mostre que A tem impureza ponderada
        menor que B se e somente se $I(1/3)/I(6/11)<11/12$. Calcule essa razão para as
        três medidas e explique que propriedade da curva da entropia a leva a escolher B.
\end{itens}

\begin{solucao}
\textbf{(a)} Na divisão A, as proporções da espécie $1$ são $1/3$ e $2/3$, cada uma num filho
com metade das mudas. Na divisão B, são $0$, numa muda, e $6/11$, nas outras $11$. Pelo Gini,
\[
  \text{A: }\ \tfrac12\cdot2\cdot\tfrac13\cdot\tfrac23+\tfrac12\cdot2\cdot\tfrac23\cdot\tfrac13=\tfrac49\approx0{,}444,
  \qquad
  \text{B: }\ \tfrac1{12}\cdot0+\tfrac{11}{12}\cdot2\cdot\tfrac6{11}\cdot\tfrac5{11}=\tfrac5{11}\approx0{,}455 .
\]
Pela entropia, como $D(1/3)=D(2/3)=\tfrac13\log_23+\tfrac23\log_2\tfrac32=\log_23-\tfrac23$,
\begin{align*}
  \text{A: }\ &\log_23-\tfrac23\approx0{,}918,\\
  \text{B: }\ &\tfrac{11}{12}D\big(\tfrac6{11}\big)
  =\tfrac1{12}\Big(6\log_2\tfrac{11}6+5\log_2\tfrac{11}5\Big)
  =\tfrac1{12}\big(11\log_211-6\log_26-5\log_25\big)\approx0{,}911 ,
\end{align*}
com $\log_26=1+\log_23$. Pelo erro de classificação, A erra $2$ mudas em cada filho, $4/12=1/3$,
e B erra as $5$ mudas da espécie $0$ no filho grande, $5/12\approx0{,}417$. O Gini e o erro
escolhem A; a entropia escolhe B. O nó tem $G=0{,}5$ e $D=1$, e as reduções correspondentes são
$0{,}056$ (A) e $0{,}045$ (B) no Gini, e $0{,}082$ (A) e $0{,}089$ (B) na entropia.

\smallskip
\textbf{(b)} Pela simetria, a impureza ponderada de A é
$\tfrac12I(1/3)+\tfrac12I(2/3)=I(1/3)$, e a de B é
$\tfrac1{12}I(0)+\tfrac{11}{12}I(6/11)=\tfrac{11}{12}I(6/11)$. Como $I(6/11)>0$,
\[
  I\big(\tfrac13\big)<\tfrac{11}{12}\,I\big(\tfrac6{11}\big)
  \iff \frac{I(1/3)}{I(6/11)}<\frac{11}{12}\approx0{,}917 .
\]
As razões são $(4/9)/(60/121)=121/135\approx0{,}896$ para o Gini,
$0{,}918/0{,}994\approx0{,}924$ para a entropia e $(1/3)/(5/11)=11/15\approx0{,}733$ para o erro: só
a entropia passa de $11/12$.

A comparação tem uma leitura simples. A divisão A leva os dois filhos da proporção $1/2$ à
proporção $1/3$; a divisão B torna pura $1/12$ das mudas e deixa as outras numa região quase tão
misturada quanto o nó, porque $6/11\approx0{,}545$ está perto de $1/2$. Com
$I(6/11)\approx I(1/2)$, a pergunta é, em essência, se $I(1/3)$ fica abaixo de $11/12$ do
máximo: $G(1/3)/G(1/2)=8/9\approx0{,}889$, mas $D(1/3)/D(1/2)\approx0{,}918$. A entropia é mais
achatada perto do máximo: trata uma região com um terço de uma espécie como quase tão impura
quanto uma equilibrada, e por isso dá pouco valor à purificação parcial de A e, em comparação,
mais valor à purificação completa de um pedaço pequeno, que é o que B faz.
\end{solucao}

\end{exercicio}

\end{document}
